\documentclass[12pt]{article}
\usepackage[utf8]{inputenc}
\usepackage{geometry}
\geometry{a4paper, margin=1.2in}
\usepackage{enumitem}
\usepackage{titlesec}
%\usepackage{parskip}
\usepackage{xcolor}
\usepackage{tfrupee}



% Formatting titles for clear readability
\titleformat{\section}{\large\bfseries\color{blue!40!black}}{}{0em}{}[\vspace{-0.5em}\rule{\textwidth}{0.5pt}]

\title{\textbf{\LARGE The Scorekeepers of Our Country:\\ \large What Do Statistics Departments Actually Do?}}
\author{Augustin Thomas}
\date{}
\begin{document}
 
\maketitle
 
\section*{Introduction}
Every state and the central government need reliable numbers to plan for the future. This is the job of the statistics departments. In Kerala, this work is done by the Department of Economics and Statistics. At the national level, it is done by the Ministry of Statistics and Programme Implementation (MoSPI). Both collect, check, and share data that helps the government make good decisions.
 
\section*{Kerala's Department of Economics and Statistics}
This department works under the Planning and Economic Affairs Department of the Kerala government. Its office is in Thiruvananthapuram, and it has branches in all 14 districts, each led by a Deputy Director. Below the district offices, there are Taluk Statistical Offices spread across the state.
 
The department started from a small survey on farm statistics back in 1949 and has grown over the decades into the main statistical body of the state. Today, its main jobs include:
 
\begin{enumerate}[leftmargin=*]
    \item \textbf{Collecting data} on agriculture, industry, prices, employment, and household income across Kerala.
    \item \textbf{Calculating Kerala's state income} (State Domestic Product) and per capita income every year.
    \item \textbf{Conducting surveys}, such as sample surveys and economic censuses, to understand how people live and work.
    \item \textbf{Advising the government} on economic and statistical matters, helping in policy and planning decisions.
    \item \textbf{Supporting local governments}, including District Planning Committees, and coordinating schemes like MPLAD funds.
    \item \textbf{Publishing reports} so that researchers, students, and the public can use accurate data about Kerala's economy.
    \item \textbf{Staying connected} with the central government's statistics office, so state and national data match up properly.
\end{enumerate}
 
In short, this department acts like Kerala's data bank, gathering facts about the economy and passing them on to planners and policymakers.
 
\section*{Government of India's Statistics Department (MoSPI)}
At the national level, statistics work is handled by the Ministry of Statistics and Programme Implementation, formed in 1999. This ministry has two parts:
 
\subsection*{1. The Statistics Wing}
Now called the National Statistical Office (NSO). It was created in 2019 by merging two older bodies, the Central Statistical Office and the National Sample Survey Office. The NSO:
\begin{itemize}[leftmargin=*]
    \item Sets the rules and standards for collecting and processing statistics across the country.
    \item Conducts big national surveys, like the Periodic Labour Force Survey and Household Consumption Expenditure Survey.
    \item Calculates the country's GDP and other national income figures.
    \item Coordinates with state statistics departments, including Kerala's, so all data follows the same methods.
\end{itemize}
 
\subsection*{2. The Programme Implementation Wing}
This wing keeps track of how government projects are progressing. It monitors big infrastructure projects, especially those costing more than \rupee150 crore, watches for delays or cost overruns, and manages the Members of Parliament Local Area Development Scheme.
 
There is also the National Statistical Commission, which oversees data quality, and the Indian Statistical Institute, a respected research and training institute.
 
\section*{Conclusion}
Both the Kerala and central statistics departments do similar work at different levels x collecting facts, running surveys, and turning raw numbers into information the government can use. Kerala's department focuses on the state's economy, while MoSPI looks at the whole country and also tracks how well government projects are being carried out. Together, they form the backbone of evidence-based planning in India.
 
\end{document}
